\documentclass[11pt,letterpaper]{article}
\usepackage[margin=1in,headheight=15pt]{geometry}
\usepackage[T1]{fontenc}
\usepackage{lmodern}
\usepackage{microtype}
\usepackage{amsmath,amssymb}
\usepackage{enumitem}
\usepackage{fancyhdr}
\usepackage{needspace}
\usepackage[hidelinks]{hyperref}
\usepackage{xcolor}
\usepackage{natbib}
\setlength{\parindent}{0pt}
\setlength{\parskip}{6pt}
\linespread{1.10}
\setlist[itemize]{leftmargin=1.3em,itemsep=3pt,topsep=4pt}
\pagestyle{fancy}
\fancyhf{}
\fancyhead[L]{\small Response to Editors and Reviewers}
\fancyhead[R]{\small MS-SPI-2025-02664}
\fancyfoot[C]{\thepage}
\renewcommand{\headrulewidth}{0.3pt}
\hypersetup{pdftitle={Response to Editors and Reviewers: MS-SPI-2025-02664},pdfauthor={}}

\newcommand{\point}[3]{%
  \Needspace{7\baselineskip}
  \subsection*{#1. #2}\phantomsection\label{pt:#1}%
  \begin{quote}\small\textbf{Comment (restated).} \textit{#3}\end{quote}%
  \textbf{Response.}\enspace\ignorespaces}
\newcommand{\locations}[1]{\par{\small\textbf{Manuscript locations:} #1.}\par}
\newcommand{\replyref}[1]{\hyperref[pt:#1]{Response #1}}

\begin{document}

\begin{center}
{\Large\bfseries Response to Editors and Reviewers}\par
\vspace{5pt}
{\large Decentralized Crypto Governance?}\par
Transparency and Concentration in Ethereum Decision-Making\par
\vspace{4pt}
Original manuscript: MS-SPI-2025-02664\par
Resubmission to the regular issue of \textit{Management Science}\par
\end{center}

Dear Professor Cong, Associate Editor, and Reviewers:

Thank you for the detailed reports and for the opportunity to submit a substantially revised paper for consideration as a new regular-issue submission to \textit{Management Science}. We understood the central message of the decision to be that the paper needed a clearer economic question, evidence more closely tied to that question, and a more precise contribution relative to existing work. We have used that assessment to guide a complete rewrite, grounding the paper in corporate finance theory, and adding new empirical analysis of contributor incentives. We appreciate the editorial team’s constructive suggestions, which have greatly improved the paper.


The main changes are:
\begin{itemize}

\item \textbf{New data linkage and analysis of contributor incentives.} This is the largest addition to the paper. We examine whether publicly observable contributions are followed by labor-market rewards. This question connects the institutional setting to the economics of open-source participation: technical work can contribute to a public project while also revealing information about the person who produces it. We thus analyze the labor market outcomes of authors after an EIP proposal. To address endogeneity and selection concerns, we match EIP authors with similar EIP authors that did not propose an EIP around that period. We construct matched cohorts, estimate static and dynamic stacked difference-in-differences models, and examine junior and senior contributors separately. We find that contributing to the Ethereum governance has important effects on the career trajectory of authors: Both role seniority and compensation rise significantly after the proposal of an EIP. 

\item \textbf{Rewritten economic motivation and literature positioning.} The abstract, introduction, literature discussion, and conclusion have been completely rewritten, connecting now  information, contributor incentives, and effective influence. The discussion distinguishes Ethereum's off-chain protocol governance from token-voting decentralized autonomous organizations (DAOs) and links the new analysis to the open-source and signaling literatures.

\item \textbf{Removal of several subsidiary empirical claims.} In an attempt to focus the paper on a single core message, we have removed the forum-engagement and proposal-success regressions, the associated company-affiliation success analysis, and the event study of EIP-document commits around developer calls. The empirical argument therefore places less weight on explaining governance outcomes through endogenous popularity and engagement measures.
\item \textbf{Reorganization and more explicit reporting.} A dedicated data subsection precedes the empirical analysis. The impact of governance on Ether prices now appears in Section II, followed by the new career analysis. New Table 2 reports the cumulative abnormal returns underlying the price figure, and the text explicitly explains the common trading-day calendar used across assets. The institutional sections and conclusion connect the evidence across governance stages.
\end{itemize}

The reports also made clear that the interpretation must match what each comparison can establish. In the responses below, we distinguish the new evidence on contributor outcomes from selection into proposals, and observed concentration from claims about authority or welfare. Where a requested analysis is not reported in the current revision, we say so explicitly.

Each comment is restated, rather than quoted verbatim, followed by our response and manuscript locations. Referee 1's five major and eleven minor comments retain their original numbering within those groups. We also address the questionnaire responses and Referee 2's report in the decision letter; overlapping Department Editor and Associate Editor concerns are grouped. Page references use the \emph{printed pages of the revised manuscript}. References to the original manuscript are identified explicitly.

\section*{Responses to the Department Editor and Associate Editor}

\point{E1}{Economic question and conceptual framework}{The Department Editor and Associate Editor find that the original paper lacks a sufficiently clear theoretical framework or research question. The evidence remains largely descriptive and insufficiently connected to governance, collective action, and information economics. A viable resubmission requires substantial reconceptualization.}

The reports prompted us to reconsider the economic argument and the evidence needed to support it. We have completely rewritten the introduction and conclusion, and grounded the paper on a framework centered around three fundamental governance issues: (i) incentives and agency problems, (ii) asymmetric information and insider trading, and (iii) effective control. The starting point is Ethereum's combination of open-source production and a tradable token: those who contribute technical work, those who implement changes, and those financially exposed to the resulting decisions need not be the same people, and thus governance rights and responsibilities are spread out along the governance value chain. 

In the new version of the paper, we thus ask three questions related to incentive alignment, insider information, and control in the Ethereum ecosystem: Does governance matter for Ether price? Do the available incentive mechanisms, like labor-market returns and company
affiliation, substitute adequately for contractual alignment? And who has the most influence in controlling the governance decision making when no formal organization exists? 

We offer four empirical tests to answer these questions: First, we show that as a proposal moves towards finalization and implementation, Ether price rises with it. Because the governance process is largely publicly visible through discussion forums and calls, this manifests as a run-up in price before the finalization of an EIP.  

Second, we show a new analysis of whether publicly observable technical contributions are associated with subsequent career gains. Section III develops the signaling argument using \citet{Spence1973} and \citet{LernerTirole2002,LernerTirole2005}, links proposal authors to employment histories, and compares career trajectories around EIP authorship. We find that role seniority and compensation increase after an EIP proposal is announced. The economic connection is that an attributable contribution to a public project may offer a private return through professional recognition. The new evidence is consistent with that channel, subject to the selection and timing qualifications in \replyref{E2}.

This mechanism gives a clearer role to the institutional evidence. Public records make contributions observable, technical complexity makes expertise relevant to producing and assessing them, and the distinction between authorship, implementation, and adoption identifies where different participants enter the process. The revision offers a conceptual framework and an empirical examination of one incentive channel. It does not establish an equilibrium account of participation, information acquisition, or the optimal distribution of control. 

Third, we document the extent of asymmetric information in the governance process. We find that governance is fully transparent: each stage, from forum ideation through proposal submission, community and developer discussion, approval, and client implementation, is publicly and immediately visible, and 89\% of participants use their real identities. But proposals average a college reading level, so this information is public yet costly to process.

Finally, we map how concentrated the control of each step of the governance process is. Ethereum’s governance concentrates and disperses at different stages. At the proposal stage, more than 600 people have put forward improvements, yet ten individuals account for 62\% of all core protocol upgrades, and the Ethereum Foundation is central, employing more authors than any other organization and affiliated with over 20\% of all EIPs. Discussion is far broader: over 1,370 people comment on proposals in online forums, and 387 participate in technical implementation calls. Implementation then narrows again, with client software development more concentrated than authorship (a Gini coefficient above 80\%) and the Foundation backing all four major client teams. Finally, adoption introduces centralization vectors that operate entirely outside the formal governance process: two stablecoin issuers account for over 90\% of stablecoin value on Ethereum, and a single oracle provider supplies the great majority of price feeds used by smart contracts, each holding an effective veto over protocol upgrades.

Overall, we believe that the new framework has significantly improved the paper. We hope you agree. 



\locations{Introduction; Section III; Section VI; Table 3; Figure 4}

\point{E2}{Selection, endogenous timing, and identification}{The editors emphasize selection into the formal EIP process, endogenous authorship, timing and content, missing forum data, and the need for a cleaner identification strategy. Proposal characteristics, approval, and market reactions cannot be treated as exogenous.}

We have changed both the empirical question and the set of claims made from the original regressions. To streamline and refocus the paper on the new framework, the forum-engagement and proposal-success regressions have been removed, including the company-affiliation success estimates. The revised paper no longer uses those specifications to explain approval through popularity or engagement. However, we introduced a new rigorous and in our view well-identified analysis of whether there is a potential return to contributing to the governance process. 

For that new question, Section III sets out an explicit comparison. A treatment event is an EIP proposal by an author with no proposal in the preceding two years. Each treated author is matched to at most three other EIP authors with no proposal anywhere in the $[-2,+2]$ event window. Matching uses seniority and log compensation in the two pre-event years. We show that treated and control groups are similar both on matched and unmatched characteristics the year before the event. The stacked specification includes author--cohort and cohort--time fixed effects, with standard errors clustered by author. Table 3 reports balance on observed pre-event characteristics and both static and dynamic estimates.

It is also important to highlight the limitations of our empirical design. A causal interpretation would require comparable changes in untreated career outcomes and an adequate account of anticipation. The design does not make submission timing random: a career transition, employer support, or work preceding formal submission could affect both authorship and subsequent outcomes. Furthermore, pre-event similarity on matched outcomes is partly imposed by construction, and while treated and control groups have similar observable characteristics before the treatment, the test cannot establish balance on unobservables. We therefore interpret the estimates as conditional evidence consistent with career returns to contribution, rather than as an isolated causal effect of signaling.

The other selection concerns remain distinct. The proposal records describe ideas that entered the formal process, the price exercise conditions on successful proposals, and the forum evidence covers identified discussion threads. Removing the original regressions reduces reliance on endogenous correlates of approval, but does not make those samples representative of unobserved proposals or discussions. We address each sample's limits in the relevant responses below.

\locations{Section I.B--C; Section II, ; Section III; Table 3}

\point{E3}{Contribution relative to existing blockchain-governance research}{The review team is not persuaded that the original descriptive findings constitute a sufficiently distinct contribution. The resubmission must explain what is new relative to existing work on concentration and governance in blockchain organizations.}

We have rewritten the literature discussion to explain the institutional comparison and the empirical addition more precisely. The relevant DAO studies examine on-chain governance through token voting, delegation, and related decisions. Ethereum's protocol governance instead proceeds through off-chain rough consensus, client implementation, and adoption. Ether ownership does not itself confer a formal protocol-governance vote. This distinction makes the incentives of technical contributors and the organization of implementation particularly relevant.

The principal empirical addition is the linkage between EIP authorship and subsequent career outcomes. It brings participant-level evidence to a mechanism emphasized in the open-source literature. The contribution we seek to establish is the empirical assessment of three pillars of corporate governance (incentives, asymmetric information, and control) across governance stages of an open-source community. The concentration statistics describe the setting in which those incentives operate. Concentration in a nominally decentralized organization is not, by itself, the claimed novelty.

\locations{Literature discussion; Section I.B; Section II; Section III; Section V}

\point{E4}{Scope and organization}{The Department Editor recommends a sharper paper built around one economic mechanism or governance friction, with one or two central claims established convincingly. The original paper tries to explain openness, concentration, proposal success, implementation, and token prices simultaneously.}

We agree with your assessment that the paper showed too many descriptive results, and few well-identified tests grounded in theory. To address this concern, we have reduced the number of empirical exercises and reorganized the paper around a more explicit economic argument. Specifically, we removed the determinants-of-engagement regressions, the proposal-success regressions and company-affiliation success analysis, and the event study of EIP-document commits around developer calls. The revised manuscript adds a dedicated data subsection and places the new career-outcome analysis immediately after the market-evidence section.

The reason for retaining the institutional sections is substantive: the career mechanism requires contributions that others can observe and attribute, while interpretation of governance outcomes requires distinguishing proposal authors from implementers and adopters. The rewritten conclusion connects those features to information, incentives, and influence. The new manuscript follows a three-part framing regarding incentives, control, and asymmetric information; the principal new empirical contribution within it is the analysis of contributor incentives.

\locations{Section I.C; Sections II--III; Sections IV--V; Section VI}

\clearpage
\section*{Response to Referees 1 and 2}

Thank you for the detailed engagement with the institutional setting and empirical analysis. Your questions about who contributes, what contributors gain, and how the descriptive evidence should be interpreted motivated the new career-outcome analysis and the rewrite of the economic argument. We address each point below.

\subsection*{Major comments}

\point{R.1}{Economic theory, costly information, and collective action}{R1: The paper needs an economic framework that guides the entire empirical analysis. Costly information acquisition, specialization, and collective action should explain who becomes informed, who invests in governance, and why others free-ride. Grossman--Stiglitz, Olson, and Ostrom are suggested as relevant foundations; adding mathematics without a coherent mechanism would not suffice. \\
R2: The paper states what it documents but does not articulate an economic puzzle, a key argument, or a theoretical debate. It needs substantial reorganization around a clear framework and well-motivated questions. The implicit tension is that token holders are economically exposed to governance despite lacking formal voting rights.}


Your comment highlights the greatest weakness of the previous version of the paper. In the revised version, we completely changed the economic framework, grounding it on both the corporate governance literature \citep[e.g.,][]{Coase1937,JensenMeckling1976}, as well as the literature on collective action (\cite{Olson1965}, \cite{Ostrom1990}), and the literature on incentives to acquire information (\cite{GrossmanStiglitz1980}). Your comment also led us to make contributor incentives a substantive empirical component of the revision. The new Section III develops a signaling mechanism: a technically demanding, publicly attributable contribution may help employers assess a contributor's skills and standing. This offers a possible private return to producing a public input, even though the protocol does not employ the contributor. The argument draws on \citet{Spence1973} and \citet{LernerTirole2002,LernerTirole2005}.

The mechanism motivates the outcomes and a comparative prediction. If contributions convey information valued in the labor market, authorship should be followed by improved career outcomes, conditional on an appropriate comparison. If the signal is more informative for less established contributors, the relationship may be stronger for those initially in junior roles. Section III examines seniority and estimated compensation, and Table 3 reports full-sample estimates and splits by pre-event seniority (levels 1--3 versus 4--7). The seniority point estimate is larger for junior contributors (0.284 versus 0.150), while the compensation estimates are close (0.093 versus 0.085 log points).

The information-cost and collective-action perspectives in your comment also clarify the role of the remaining evidence. Open access does not remove the expertise needed to assess proposals, and a possible career return helps explain why participation can be privately attractive despite the public nature of the output. The evidence does not distinguish signaling from learning or network formation, measure free-riding, or establish an efficient allocation of governance effort. Our response is therefore a conceptual mechanism linked to an empirical comparison, rather than a formal solution to the broader information-acquisition or collective-action problem.

\locations{Introduction; Section III, ; Section IV; Section VI; Table 3}

\point{R.2}{Selection into proposals and strategic participation}{R1: Only submitted EIPs are observed, although proposal formation follows feedback and depends on proposers' incentives. Explain who proposes, what returns they expect, and whether those returns encourage continued development. Timing, content, co-authorship, and discussion participation may be strategic. Concentration may reflect efficient specialization. Stagnation can reflect learning about proposal quality, so stagnant and withdrawn proposals need distinct economic interpretations. Any persuasion framework should be coherent with the broader theory.\\
R2: Methodological concern 1: only formal EIPs are observed. Author characteristics, social influence, and engagement may be endogenous to quality, and timing may be strategic. Address selection through an appropriate design or at least acknowledge the limitation. Changes in selection also complicate concentration trends.}

We agree that a submitted EIP is the outcome of earlier choices and feedback. Section I.B describes discussion before formal submission: for example, Ethereum Research (ethresear.ch) is a discussion forum where ideas about future protocol improvement are discussed. However, measuring the link between research ideas and formalized EIPs is very challenging, and very noisy at best. The authorship distribution is therefore conditional on entry into the EIP process. Strategic choice of timing, content, and collaborators remains relevant, as does the possibility that frequent contributors are selected for ability, resources, or expected benefits.

The principal addition addresses your questions about who contributes and what contributors may gain. We linked 469 of the 685 EIP authors to employment histories and constructed a matched comparison around proposal authorship. Table 3, Panel A, describes that career sample: 64\% of treated observations are classified in engineering occupations and 91\% have an identified employer in the pre-event year. These are characteristics of the matched sample, not estimates for all potential contributors.

In the full matched estimation sample, Table 3 reports relative post-event changes of 0.221 points on the seven-level seniority scale (standard error 0.083) and 0.090 log points in estimated total compensation (standard error 0.039). The latter corresponds to approximately 9\%. The post period includes the proposal year and the following two years. The comparison is with other EIP authors without a proposal in the same event window, not with individuals who never participate. Its construction and identifying limits are described in \replyref{E2}.

The positive estimates are consistent with career benefits, but require care in interpretation. Compensation is estimated by Revelio Labs from employment characteristics and external surveys, rather than observed pay, and seniority is one input to that prediction. The two outcomes therefore do not provide complete independent evidence of a pay and promotion mechanism. 

In order to streamline and align the empirical analysis with the new framework, we have also removed the proposal-success regressions that classified stagnant and withdrawn EIPs together as failures. The institutional discussion retains the relevant distinction: stagnation follows a period without progress and can be reversed, whereas withdrawal is irreversible. Learning about quality or the option to resume work can help interpret that distinction, but neither explanation is identified by the status records alone. Likewise, specialized skills and organizational support can generate concentrated contribution without implying capture. A structural account of proposal formation or Bayesian persuasion is outside the evidence established here.

\locations{Section I.B; Section III; Section V.A; Section V.E; Table 3}

\point{R.3}{Related literature and network versus firm value}{R1: Engage substantively with DAO-governance research, including Appel and Grennan, Cong et al., and Fan, Shu and Xie. Explain the incremental insight, the connection to optimal governance and value, and how network value differs from firm value.\\
R2: Existing studies already document concentrated blockchain governance. Explain the additional insight and whether observed authorship concentration and increasing network clustering represent efficient specialization or problematic capture.}

We have rewritten the literature discussion to make the comparison with the studies you identify more explicit. The revised discussion relates our setting to \citet{AppelGrennan2023,Appel-Grennan2024}, \citet{Cong-et-al2025}, \citet{Fan-et-al2025}, and \citet{Bongaerts-et-al2025}. Ethereum's protocol-level rough-consensus process separates token ownership from formal decision rights, whereas the cited DAO literature examines voting, delegation, and governance within token-based organizations. Ethereum's governance more closely resembles a traditional open-source community with indirect links to token ownership, while the DAOs studied in this literature resemble a corporate governance structure in which tokenholders vote to approve changes proposed by a management team or by community members. Off-chain rough consensus avoids the concentration risk inherent in token-weighted voting: as \citet{AppelGrennan2023,Appel-Grennan2024} document, a small number of entities can exert outsized control over DAO voting outcomes, whereas Ethereum's process requires buy-in across client teams and the broader developer community rather than a majority vote among tokenholders. This paper therefore focuses on an off-chain governance process, with incentive structures and goals distinct from those of on-chain DAO governance.

The new analysis on labor market consequences, paired with the impact of protocol changes on Ether price, closes the loop between rough-consensus off-chain governance, the incentives of contributors, and network value maximization. It asks what rewards may accrue to people undertaking visible technical work, drawing on the open-source literature and the added \citet{Spence1973} and \citet{LernerTirole2002} references. The intended incremental insight is about contributor incentives in this governance setting, rather than another demonstration that blockchain participation is concentrated.

\locations{Literature discussion; Section II; Section III; references}

\point{R.4}{Price run-up, market efficiency, and counterfactual}{R1: Explain what the reported 12\% increase says about information aggregation or market efficiency, whether it is efficient, and what counterfactual is appropriate. A transparent process reveals information gradually, making a conventional event study problematic. Improve identification in a way aligned with the framework.\\
R2: Price movements may reflect anticipation or selection rather than governance quality. A run-up does not establish causality or show that token holders benefit from the existing governance process.}

We agree with the referee that a conventional event study, one that treats the pre-event window as a clean untreated pre-period and looks for a sharp reaction at the event date, is not well suited to a setting where a transparent process reveals information gradually. Our design instead treats the pre-event window itself as the object of interest: the hypothesis we test is that information arrives gradually rather than as a single surprise at finalization.

We take all successfully finalized Core EIPs and assume that markets are efficient, so that as an EIP moves through the governance process, information about its eventual approval is released and the market updates its assessment of the finalization probability. If EIPs are relevant for Ether's value, we should expect price to rise gradually over the run-up to finalization. The gradual shape of the run-up we document is therefore consistent with markets being efficient, with information being released transparently throughout the governance process, and with EIPs affecting token value.

The revision makes this reasoning prominent in the introduction and places the price analysis in Section II. 

We have added Table 2, which reports the cumulative abnormal returns underlying Figure 3 at selected horizons, and made the common trading-day calendar explicit. Ether, Bitcoin, and the S\&P 500 use the same calendar, with weekend and holiday price changes included in the following trading-day return. These additions make the reported path easier to inspect. 

The appropriate interpretation is a price pattern conditional on selected governance milestones. Bitcoin provides a crypto-market reference, and the S\&P 500 provides a broad equity-market reference and is subject to substantial sample-period drift relative to Ether. Neither supplies the counterfactual Ether price absent the proposal. 

The observed run-up is compatible with gradual updating about expected protocol changes, but compatibility is not a test of efficient information aggregation. The exercise cannot determine whether prices incorporate information correctly, isolate the value created by approval, or rank governance arrangements by welfare. We retain it as market evidence relevant to the economic stakes of governance, with this narrower interpretation.

\locations{Introduction; Section II; Figure 3; Table 2}

\point{R.5}{Network measures and their economic interpretation}{Justify betweenness as a measure of information bridging, consider eigenvector centrality for reputation and consensus formation, investigate communities or factions and their relationship to polarization, and provide statistical tests for network properties. The model should also guide whether other node characteristics belong in regressions.}

We agree that the network statistic should be matched to the economic claim. In the retained graph, nodes are authors and links record co-authorship. Betweenness describes positions lying on shortest paths between collaborators. In the network analysis literature it is commonly used to proxy for either information flow or brokering positioning. On the other hand, as you correctly stated, eigenvector centrality addresses connectedness to well-connected collaborators and is commonly considered a proxy for influence and consensus building (see e.g. \cite{WassermanFaust1994}). Neither statistic alone identifies governance authority.

As with other tests, the revision removes the engagement and proposal-success regressions that used network position, so the paper no longer relies on betweenness or the associated combinations of node characteristics to explain approval. The retained graph describes collaboration, while the contribution rankings and concentration measures describe the allocation of recorded activity, but the new version of the paper no longer mentions any measure of network centrality.


\locations{Section V.A; Figure 9; Figure 10; Figures 11--13}

\subsection*{Minor comments}

\point{R.m1}{Partial-year 2023 observations}{Incomplete 2023 data can create artificial endpoint effects. Extend the collection or exclude 2023.}

We agree that some of the graphs might suggest a decline in activity simply because the data collection ends in the middle of 2023. It does not affect any of the analysis, but it might be visually misleading. To address this concern, we decided to add an asterisk next to the 2023 year in all graphs, and a note in the caption, to clarify that the data only goes up to June 2023. However, if you feel strongly about it, we are happy to exclude 2023.

\locations{Section I.C; Figures 2, 12--15, and 21}

\point{R.m2}{Excluded Meta and Informational EIPs}{Explain the purpose of Meta and Informational EIPs, justify their exclusion, and report their percentage of total EIPs.}

The intended sample is proposals for technical changes and application standards. Meta proposals concern processes, including governance or development procedures; Informational proposals provide design information or guidance without proposing a new feature.\footnote{See \href{https://eips.ethereum.org/EIPS/eip-1}{EIP 1: EIP Purpose and Guidelines}, under EIP Types.} They are excluded under the sample definition stated in Section I.C, footnote 8. Overall, 639 EIPs were proposed during our sample period; 27 of them (4.2\%) were either Meta (21) or Informational (6). We added these statistics in the paper. 



\locations{Section I.B; Section I.C and footnote 8}

\point{R.m3}{Missing FEM discussion data}{Only 346 of 612 EIPs have FEM discussion data. Analyze missingness and whether these proposals are systematically different, including whether discussion is selected by the uncertainty of approval.\\
R2: Methodological concern 2: 346 of 612 proposals have forum data, leaving approximately 44\% without coverage. Test whether missingness is related to proposal characteristics or whether technical proposals bypass the forums; otherwise engagement findings may describe selected proposal types.}

The original count of 346 omitted threads in several forum categories. In the new draft we corrected this omission and find that 448 of the 612 EIPs (73\%) are discussed on FEM. If you include discussions on GitHub and Ethereum Research, collected in the new draft, the total EIPs with discussion in at least one open venue reaches 588 of 612 EIPs (96\%). Missing EIPs are due to the fact that the FEM site launched on February 16, 2018, and thus early EIPs had no formal venue for discussion.

In the revised version of the paper, we have removed the engagement regressions and the proposal-success regressions that used the community-engagement index. This removes those specific inferential uses of the incomplete forum sample. The remaining statistics describe observed FEM activity. In the current revision, FEM data are used only to describe participation; they are not an outcome or regressor in any analysis. The broad-participation conclusion does not rest on FEM alone. It uses FEM, GitHub, and AllCoreDevs calls.



\locations{Section IV; Section V.B; Figure 14}

\point{R.m4}{Cluttered network visualization}{Improve the layout or highlight key nodes. If centrality is the purpose, a table may be preferable. The visualization alone cannot establish the presence or absence of communities.}

Thank you for pointing out the lack of clear visualization of the nodes and clusters. We decided to use a different algorithm to plot distances between nodes, to make the layout and the presence of communities cleaner. In the new graph, we position nodes using the Fruchterman-Reingold force-directed algorithm, with attraction between nodes weighted by the number of EIPs co-authored. Node color indicates community membership, identified via greedy modularity maximization on the weighted co-authorship network (12 communities). 

Figures 9 and 11 separately report author rankings and concentration, which are more useful for quantitative comparisons of contribution activity.


\locations{Figures 9--11; Section V.A}

\point{R.m5}{Readability and adoption}{The readability analysis is interesting but underdeveloped. How does readability correlate with adoption? This may be a separate research question.}

We appreciate the distinction in your comment between an interesting extension and what the present paper needs to establish. The retained readability evidence describes the written material that participants encounter. In the current version of the paper, we argue that transparency has two components, availability (covered first, via the forum/venue discussion) and comprehensibility (covered here, via Flesch Reading Ease + length). The 39.1 score and the self-referential comparison to the paper’s own readability (43.5) are used to argue that EIPs demand a college-level reading ability, i.e., information is public but costly to process.

 The original proposal-success regressions have been removed, so the revised argument no longer relies on readability as an explanation of approval.

\locations{Section IV; Figure 6; Table 1}

\point{R.m6}{Theoretical justification of controls}{Explain the economic logic behind the controls in the original Table 2 regression specifications.}

We have removed the original Table 2 engagement regressions and Table 3 proposal-success regressions, including their combinations of readability, social-following, and network variables. The revised paper does not rely on those specifications.

The new career analysis has a more direct link between question and specification. Pre-event seniority and compensation guide matching because subsequent changes in those outcomes are the object of the comparison. Author--cohort fixed effects absorb stable differences within a cohort, and cohort--time fixed effects absorb common time variation. Other pre-event characteristics are reported for balance assessment. These choices make the comparison explicit, while leaving the selection and timing qualifications in \replyref{E2}.

\locations{Section III and equation (1); Table 3}

\point{R.m7}{Regulatory motivation}{R1: The regulatory motivation is interesting but gets lost and is not revisited after the descriptive analysis.\\
R2: The paper does not explain whether developers should worry about concentration, whether the Foundation should reduce or increase its role, or whether stablecoin issuers are systemically important. Transparency does not necessarily substitute for disclosure rules, and concentration does not necessarily justify securities-style regulation.}

We have rewritten the introduction and conclusion and added synthesis at the ends of the transparency and concentration sections. These passages connect public access, technical information-processing demands, contributor incentives, and influence across governance stages. The discussion of the Foundation and infrastructure providers explains why organizational affiliations and dependencies remain relevant even when formal protocol voting rights are absent.

The implication we draw is about what a regulator needs to examine to assess a protocol's governance: access to information, the ability to interpret it, contributor incentives and affiliations, and dependencies affecting implementation or adoption. The evidence does not determine a legal classification, measure unlawful trading, or establish that a particular disclosure regime is optimal. This keeps the regulatory motivation connected to the empirical setting without turning empirical evidence into legal conclusions.

\locations{Introduction; end of Section IV; end of Section V; Section VI}

\point{R.m8}{Interpretability and marginal effects}{Report marginal effects to make the original proposal-success coefficients easier to interpret.}

Thank you for emphasizing interpretation in economically meaningful units. The proposal-success model has been removed, so there is no retained specification for which to add the requested marginal effects.

For the new analysis, the full-sample estimates in Table 3 are 0.221 seniority-scale points and 0.090 log points in estimated total compensation. We describe the latter as approximately 9\% and the former in points on the seven-level scale. This avoids interpreting a change in an ordinal seniority score as a conventional percentage promotion effect.

\locations{Section III; Table 3, Panels B--C}

\point{R.m9}{Validity of the readability measure}{Flesch Reading Ease may be inappropriate for technical documents. Consider alternative measures, examples, or robustness checks.}

We agree that Flesch Reading Ease was designed for general prose and does not capture the domain-specific terminology that makes a technical document hard to read. Flesch Reading Ease penalizes syllable/character count, which punishes legitimate technical vocabulary like “cryptographic” or “asynchronous” rather than genuinely convoluted writing. Unfortunately, alternative formulas (Flesch-Kincaid Grade Level, Gunning Fog, SMOG, Coleman-Liau, ARI) all use the same syllable-or-character-count logic. 

 We note, however, that it has been applied to technical writing elsewhere: \citet{VainshteinEtAl2019} compute it for a corpus of nearly 30,000 computer science conference papers and report a mean of 40.12, comparable to the mean of 39.1 we report for EIPs. In the revised paper, we now comment on how EIP readability compares with computer science papers. 


\locations{Section IV; Figure 6; Table 1}

\point{R.m10}{Substantive versus formatting commits}{Use commit messages to distinguish substantive changes from minor formatting edits.}

We agree that commit counts cannot distinguish substantive work from mechanical edits. For EIP documents, the manuscript excludes bots and distinguishes author from outside contributions. Those filters identify contributors. They do not classify the technical importance of an edit.

To address your concern, we expanded the commit data collection, and downloaded GitHub pull-request/issue comments and formal review submissions for each EIP, reported in the new Panel B of Table 4. On average, an EIP receives about 14 comments and 18 reviews, compared with about 7 commits. A review is a closer proxy for substantive engagement than a commit: posting one requires reading and responding to the technical content of a specific proposal, rather than making or accepting a text change, and unlike commits it is not dominated by authors. This addition does not change our broad-participation conclusion. It reinforces it. Average GitHub review activity (18 per EIP) exceeds average FEM forum comments (15 per EIP), indicating that code-level review is at least as active a venue for external scrutiny of EIPs as the open discussion forum.

We have removed the developer-call commit event study and its claim about a tripling of update frequency around calls. The remaining commit distributions describe recorded repository activity. 



\locations{Section IV; Section V.C; Figures 5 and 16}







\point{R.m11}{Theory-guided empirical choices}{Methodological concern 3: a high-level framework should explain efficient specialization versus harmful concentration, the role of transparency, the outcomes of interest, and the control variables. The original specifications, including readability controls, appear arbitrary.}

The original regressions containing the criticized readability and engagement controls have been removed, so the objection to those specification choices is no longer relevant. Readability remains a descriptive text measure.

The broader request for economic motivation guides the revised framing described in \replyref{R2.1}. New Section III links observable technical contributions to potential career returns, motivating seniority and estimated compensation as outcomes and matching on pre-event career characteristics. The institutional evidence connects information access and contributor incentives to implementation and adoption dependencies. This provides a clearer rationale for the empirical choices, while the distinction between efficient specialization and capture remains subject to the limits set out in \replyref{R2.2}.


\bigskip
We are grateful for the review team's close reading and guidance. The new employment linkage and analysis, rewritten economic argument, and removal of several subsidiary specifications reflect the substantive work prompted by the reports. We appreciate the opportunity to have the revised work considered as a new submission.


\newpage
\bibliographystyle{aea}
\bibliography{bib.bib}


\end{document}
